\documentclass[12pt,a4paper,oneside]{article}

\usepackage[QX]{polski}

\usepackage[utf8]{inputenc}
\usepackage{latexsym}
\usepackage{tgpagella}
\usepackage{lmodern}
\usepackage{amsmath,amsthm,amsfonts,amssymb,alltt}
\usepackage{epsfig}
\usepackage{pdflscape}
\usepackage{caption}
\usepackage{indentfirst}
\usepackage{float}
%\usepackage{showkeys}
\bibliographystyle{plabbrv}


\usepackage{color}
\usepackage[polish]{babel}
\usepackage{datetime2}
\usepackage[x11names,dvipsnames,table]{xcolor}
\usepackage{hyperref}
\hypersetup{
pdfauthor={Roman Czapla, Olaf Bar},
colorlinks=True,
linkcolor=darkgray,  % color of internal links (change box color with linkbordercolor)
citecolor=BrickRed,  % color of links to bibliography
filecolor=Magenta,   % color of file links
urlcolor=BlueViolet}	%%pdfpagemode=FullScreen}

% diagramy, grafy itp.
\usepackage{tikz}
\usetikzlibrary{positioning}
\usetikzlibrary{arrows}
\usetikzlibrary{arrows.meta}
\usetikzlibrary{chains,fit,shapes,calc}
\tikzset{main node/.style={circle,fill=blue!20,draw,minimum size=1cm,inner sep=0pt}}

% algorytmy
\usepackage[linesnumbered,lined,commentsnumbered]{algorithm2e}
\SetKwFor{ForEach}{for each}{do}{end for}%
\SetKwFor{ForAll}{for all}{do}{end for}%
\newenvironment{myalgorithm}
{\rule{\textwidth}{0.5mm}\\\SetAlCapSty{}\SetAlgoNoEnd\SetAlgoNoLine\begin{algorithm}}{\end{algorithm}\rule{\textwidth}{0.5mm}}


%---------------------
\overfullrule=2mm
\pagestyle{plain}
\textwidth=15cm \textheight=685pt \topmargin=-25pt \linespread{1.3} 
\setlength{\parskip}{0pt}
\setlength\arraycolsep{2pt}
\oddsidemargin = 0.9cm
\evensidemargin =-0.1cm

\captionsetup{width=.95\linewidth, justification=centering}
%---------------------




\newtheorem{tw}{Twierdzenie}[section]
\newtheorem{lem}[tw]{Lemat}
\newtheorem{co}[tw]{Wniosek}
\newtheorem{prop}[tw]{Stwierdzenie}
\theoremstyle{definition}
\newtheorem{ex}{Przykład}
\newtheorem{re}[tw]{Uwaga}
\newtheorem{de}{Definicja}[section]



\newcommand{\bC}{{\mathbb C}}
\newcommand{\bR}{{\mathbb R}}
\newcommand{\bZ}{{\mathbb Z}}
\newcommand{\bQ}{{\mathbb Q}}
\newcommand{\bN}{{\mathbb N}}
\newcommand{\captionT}[1]{\caption{\textsc{\footnotesize{#1}}}}
\renewcommand\figurename{Rys.}

\numberwithin{equation}{section}
\renewcommand{\thefootnote}{\arabic{footnote})}
%\renewcommand{\thefootnote}{\alph{footnote})}



\begin{document}

% --------------------------------------------
% Strona tytułowa
% --------------------------------------------

\thispagestyle{empty}
\begin{titlepage}
\begin{center}\Large
Uniwersytet Komisji Edukacji Narodowej w Krakowie\\
\large
Instytut Bezpieczeństwa i Informatyki\\
\vskip 10pt
\end{center}
\begin{center}
\centering \includegraphics[width=1.0\columnwidth]{images/logo.png}
\end{center}

\begin{center}
 {\bf \fontsize{14pt}{14pt}\selectfont PROJEKT INŻYNIERSKI \\ DOKUMENTACJA UŻYTKOWA}
\end{center}
\vskip 5pt
\begin{center}
 {\bf \fontsize{22pt}{22pt}\selectfont System WasteFree Cloud do efektywnego zarządzania odpadami}
\end{center}

\begin{center}
 {\fontsize{12pt}{12pt}\selectfont wykonany przez: }
\end{center}
\begin{center}
 {\bf\fontsize{16pt}{16pt}\selectfont Dawid Plata}\\
 {\fontsize{12pt}{12pt}\selectfont Nr albumu: 165303 \\\&\\}
 {\bf\fontsize{16pt}{16pt}\selectfont Michał Szymbara}\\
 {\fontsize{12pt}{12pt}\selectfont Nr albumu: 165297 \\\&\\}
 {\bf\fontsize{16pt}{16pt}\selectfont Michał Wechmann}\\
 {\fontsize{12pt}{12pt}\selectfont Nr albumu: 165298}
\end{center}
\begin{center}
 {\fontsize{12pt}{12pt}\selectfont pod opieką:}\\
 {\bf\fontsize{12pt}{12pt}\selectfont tytuł, imię i nazwisko opiekuna projektu}
\end{center}

%\mbox{}
\vspace*{\fill}
%\vskip 50pt
\begin{center}
\large
Kraków \the\year\\
(ostatnia aktualizacja: \DTMcurrenttime,\;\today)
\end{center}
\end{titlepage}
\setcounter{page}{0} 
\newpage\null\thispagestyle{empty}
%\setcounter{page}{0} 
%\newpage
%\thispagestyle{empty}

\tableofcontents


\newpage

\section{Szczegółowa dokumentacja użytkowa}

\subsection{Role użytkowników i uprawnienia}
System WasteFree Cloud przewiduje kilka ról użytkowników, które różnią się zakresem dostępnych funkcji. Uprawnienia są nadawane automatycznie na podstawie typu konta oraz autoryzacji (logowania). Dzięki temu użytkownik widzi w aplikacji jedynie te elementy, które są dla niego przeznaczone, a operacje wrażliwe (np. obsługa zleceń) są ograniczone do odpowiednich ról.

\subsubsection{Użytkownik (User)}
Rola \textbf{Użytkownik} przeznaczona jest dla osób korzystających z aplikacji w celu organizacji i rozliczania odbiorów odpadów w ramach grupy.
\begin{itemize}
\item zarządzanie własnym kontem: rejestracja, aktywacja konta, logowanie,
\item wgląd i edycja profilu (m.in. podstawowe dane, adres, miasto),
\item korzystanie z portfela: podgląd salda oraz wykonywanie transakcji (wpłata oraz wypłata),
\item tworzenie i przegląd grup odpadów oraz współpraca z innymi użytkownikami (zaproszenia),
\item składanie zamówień odbioru odpadów, kalkulacja kosztu oraz realizacja płatności,
\item korzystanie ze skrzynki odbiorczej (powiadomienia i akcje) oraz czatu grupowego.
\end{itemize}

\subsubsection{Administrator odpadów (Garbage Admin)}
Rola \textbf{Administrator odpadów} przeznaczona jest dla podmiotów obsługujących zlecenia odbioru odpadów. Użytkownik z tą rolą ma dostęp do funkcji związanych z przyjmowaniem oraz realizacją zamówień.
\begin{itemize}
\item przegląd zleceń oczekujących w przypisanym mieście oraz akceptacja zleceń,
\item wgląd w zlecenia aktualnie przypisane/aktywne,
\item uzupełnianie informacji o kosztach (np. zgłoszenie opłaty utylizacyjnej wraz z załącznikiem),
\item akceptacja wymaganej zgody (treści zgody) dla administratora odpadów.
\end{itemize}

\subsubsection{Administrator systemu (Admin)}
Rola \textbf{Administrator systemu} jest przewidziana jako rola o najszerszych uprawnieniach. W aktualnej wersji aplikacji administrator ma dostęp do wszystkich zakładek i treści w interfejsie, natomiast nie posiada dodatkowych, dedykowanych funkcji administracyjnych ponad standardowe możliwości dostępne w aplikacji.

\subsection{Nawigacja i elementy interfejsu}
Interfejs użytkownika w aplikacji WasteFree Cloud jest oparty o widoki (strony) dostępne z poziomu nawigacji oraz tzw. portalu użytkownika. Po zalogowaniu użytkownik uzyskuje dostęp do panelu \textit{Portal}, w którym znajdują się funkcje operacyjne (portfel, skrzynka odbiorcza, grupy, zamówienia, profil). Część elementów interfejsu (np. przyciski akcji) może różnić się w zależności od roli.

\subsubsection{Strona główna}
Strona główna jest punktem startowym aplikacji. Z tego miejsca użytkownik może przejść do logowania/rejestracji lub do portalu (jeżeli jest już zalogowany).

\begin{figure}[H]
\centering
\fbox{\parbox{0.9\linewidth}{\vspace{2.2cm}\centering\textit{Miejsce na zrzut ekranu: Strona główna (Home)}\vspace{2.2cm}}}
\captionT{Strona główna aplikacji}
\end{figure}

\subsubsection{Portal użytkownika}
Portal użytkownika stanowi główny obszar pracy w aplikacji. W portalu dostępne są podstrony związane z realizacją funkcjonalności systemu, m.in.:
\begin{itemize}
\item \textbf{Portfel} (saldo i transakcje),
\item \textbf{Skrzynka odbiorcza} (wiadomości oraz akcje),
\item \textbf{Zamówienia odbioru} (tworzenie i podgląd),
\item \textbf{Grupy} (lista, zarządzanie i panel grupy),
\item \textbf{Profil} (dane użytkownika).
\end{itemize}

\begin{figure}[H]
\centering
\fbox{\parbox{0.9\linewidth}{\vspace{2.2cm}\centering\textit{Miejsce na zrzut ekranu: Widok Portalu (Portal)}\vspace{2.2cm}}}
\captionT{Widok portalu użytkownika}
\end{figure}

\subsubsection{Topbar, przełącznik języka i nawigacja}
W górnej części aplikacji dostępny jest pasek nawigacyjny (topbar), który ułatwia przechodzenie pomiędzy zakładkami. W interfejsie występuje również przełącznik języka, dzięki któremu użytkownik może zmienić język wyświetlanych treści (o ile są dostępne tłumaczenia).

Elementy nawigacji mogą obejmować m.in. przejście do portfela, skrzynki odbiorczej, grup, profilu oraz innych modułów. Dodatkowo w topbarze może pojawiać się informacja o liczbie nowych powiadomień (np. w skrzynce odbiorczej).

\begin{figure}[H]
\centering
\fbox{\parbox{0.9\linewidth}{\vspace{2.2cm}\centering\textit{Miejsce na zrzut ekranu: Topbar oraz przełącznik języka}\vspace{2.2cm}}}
\captionT{Nawigacja w aplikacji (topbar i przełącznik języka)}
\end{figure}

\subsubsection{Komunikaty, walidacja i obsługa błędów}
Podczas korzystania z aplikacji użytkownik może otrzymywać komunikaty o poprawnym wykonaniu operacji (np. zapis profilu, utworzenie zamówienia) oraz komunikaty o błędach. Walidacja danych (np. w formularzach rejestracji lub płatności) ma na celu wskazanie brakujących lub niepoprawnych informacji jeszcze przed wysłaniem danych.

W przypadku wystąpienia błędu po stronie systemu aplikacja prezentuje odpowiedni komunikat, a użytkownik może ponowić operację lub skorygować dane wejściowe.

\begin{figure}[H]
\centering
\fbox{\parbox{0.9\linewidth}{\vspace{2.2cm}\centering\textit{Miejsce na zrzut ekranu: Przykładowy komunikat błędu/walidacji}\vspace{2.2cm}}}
\captionT{Przykładowe komunikaty i walidacja danych}
\end{figure}

\subsection{Konto i uwierzytelnianie}
\subsubsection{Rejestracja konta}
\subsubsection{Aktywacja konta (link e-mail)}
\subsubsection{Logowanie i wylogowanie}

\subsection{Profil użytkownika}
\subsubsection{Podgląd i edycja profilu}
\subsubsection{Adres i miasto (dostępne miasta)}
\subsubsection{Avatar użytkownika (upload)}
\subsubsection{Statystyki użytkownika}

\subsection{Portfel (Wallet)}
\subsubsection{Saldo}
\subsubsection{Dostępne metody płatności}
\subsubsection{Wpłata środków}
\subsubsection{Wypłata środków}

\subsection{Grupy odpadów (Garbage Groups)}
\subsubsection{Lista grup}
\subsubsection{Rejestracja nowej grupy}
\subsubsection{Szczegóły grupy}
\subsubsection{Zarządzanie grupą (edycja danych)}
\subsubsection{Zapraszanie użytkowników do grupy}
\subsubsection{Zaproszenia oczekujące (akceptacja/odrzucenie)}
\subsubsection{Usuwanie użytkownika z grupy}

\subsection{Panel grupy i czat (Group Chat)}
\subsubsection{Historia wiadomości}
\subsubsection{Dodawanie wiadomości}

\subsection{Zamówienia odbioru odpadów (Garbage Orders)}
\subsubsection{Utworzenie zamówienia}
\subsubsection{Kalkulacja kosztu zamówienia}
\subsubsection{Płatność za zamówienie}
\subsubsection{Moje odbiory (lista zamówień użytkownika)}
\subsubsection{Szczegóły zamówienia}
\subsubsection{Dodatkowa opłata utylizacyjna (udział w koszcie)}

\subsection{Skrzynka odbiorcza (Inbox) i powiadomienia}
\subsubsection[Licznik wiadomości i powiadomienia]{Licznik wiadomości i powiadomienia w czasie rzeczywistym}
\subsubsection{Lista wiadomości (paginacja)}
\subsubsection{Akcje na wiadomościach (akceptacja/odrzucenie)}
\subsubsection{Usuwanie wiadomości}
\subsubsection{Czyszczenie skrzynki}

\subsection{Funkcje dla Administratora odpadów (Garbage Admin)}
\subsubsection{Zgoda administratora odpadów (treść zgody)}
\subsubsection{Akceptacja zgody}
\subsubsection{Zlecenia oczekujące (waiting)}
\subsubsection{Akceptacja zlecenia}
\subsubsection{Zlecenia przypisane/aktywne (assigned/current)}
\subsubsection{Zgłoszenie opłaty utylizacyjnej (załącznik)}

\subsection{Przykładowe działanie systemu (scenariusze)}
\subsubsection{Rejestracja i aktywacja konta}
\subsubsection{Utworzenie grupy i zaproszenie użytkowników}
\subsubsection{Złożenie zamówienia oraz płatność}
\subsubsection{Komunikacja w czacie grupy}
\subsubsection[Obsługa zaproszeń i akcji]{Obsługa zaproszeń i akcji z poziomu skrzynki odbiorczej}
\subsubsection[Scenariusz Garbage Admin]{Scenariusz Garbage Admin: akceptacja zlecenia i rozliczenie opłaty}





\renewcommand\refname{Literatura (jeżeli wymagana)}
\bibliography{references}
\addcontentsline{toc}{section}{Literatura}
% --------------------------------------------------------------------
%%%%%%% odkomentować gdy bibliografia ma być wewnątrz dokumentu
% --------------------------------------------------------------------
%\begin{thebibliography}{11}
%
%\addcontentsline{toc}{section}{Literatura}
%
%\bibitem{ZAN}
%C. Zannoni and P. Pasini, 
%\emph{Advances in the Computer Simulatons of Liquid Crystals}, Kluwer Academic Publishers, 2000.
%
%\end{thebibliography}

\end{document}

